% 作者修订稿：2026-10-27。Fig4 使用 Rank、Order 标签；不修改冻结论文全文。
% Fig5 对应 vertical_legend_20261027；选点规则和完整统计见该图现有来源。
\begin{figure}[htbp]
\centering
\includegraphics[width=\textwidth]{figures/Fig4.pdf}
\caption{\textbf{患者间期时空模式约束的空间 E--I 网络构建与优化。}
\textbf{a，}局部二维兴奋--抑制网络结构。$E\to E$ 连接采用椭圆形各向异性连接核，其余局部通路采用圆形各向同性连接核；上下两层用于区分 E、I 群体，表示 b 中的局部网络。
\textbf{b，}患者电极布局下的完整空间网络。依据患者 SEEG 触点几何设置二维网络中的虚拟电极，插图显示触点邻域的局部连接与高斯采样权重。
\textbf{c，}患者间期事件约束的参数优化。模型与患者事件的秩次模式、传播顺序和触点参与分布分别构成损失 $\mathcal{L}_r$、$\mathcal{L}_o$ 和 $\mathcal{L}_p$，总损失为 $J=(\mathcal{L}_r+\mathcal{L}_o+\mathcal{L}_p)/3$；两侧显示代表事件的时频图和训练秩次，圆点表示谱图质心，短竖线表示训练秩次对应的时刻。
\textbf{d，}无 core 网络的背景活动。均匀背景输入下，网络呈稀疏放电和低幅 LFP 代理信号，虚线圆标示与 e 相同的局部读出区域，右下显示该区域内的 E、I 神经元放电。
\textbf{e，}单 core 网络的局部簇发活动。紫色填色圆区标示 core 的位置与范围，圆外保持白底。右侧显示 core 区域内 847 个 E 神经元和 191 个 I 神经元的放电及对应 LFP 代理信号，短暂同步簇发伴随明显的瞬态峰值；d、e 的 LFP 纵轴分别缩放，a.u. 表示任意单位。
\textbf{f，}core 数量对拟合表现的影响。双 core 已获得相对单 core 的主要损失下降，增加至三、四 core 的额外改善较小；曲线和阴影表示示例患者每组 4 次优化重复在共同前 31 个 epoch 内的最佳训练损失均值$\pm$标准差。
\textbf{g，}参数搜索中的损失变化。随参数搜索推进，三类损失总体向较低水平集中，但不同分量的改善程度并不一致。
\textbf{h，}低损失配置的 core 位置分布。413 个可评分配置中损失最低的 20\%（83 个）形成两个主要位置集中区；灰点表示搜索位置，轮廓包围各 core 位置密度的 90\% 质量，十字标示密度峰的平面位置。
\textbf{i，}$E\to E$ 连接方向对拟合损失的影响。彩色曲线表示三个损失分量，黑线表示总损失，阴影表示时间块重采样标准差；在测试的 13 个方向中，总损失在相对参考轴偏移 $0^\circ$ 时最低，竖虚线标示该位置。}
\label{fig4}
\end{figure}

\begin{figure}[htbp]
\centering
\includegraphics[width=\textwidth]{figures/Fig5.pdf}
\caption{\textbf{双 core 空间 E--I 网络重现患者双向间期传播模式。}
\textbf{a，}优化后代表性网络的空间结构。紫色虚线圈标示 Core A、Core B，红色椭圆及虚线标示 $E\to E$ 连接的各向异性主轴，灰色线和圆点表示虚拟电极杆及触点。
\textbf{b，}同一网络中的双向传播读出。连续 30--80\,Hz 虚拟触点活动中出现 MTA 和 MTB 事件（阴影），两类事件表现出相反的触点招募顺序。
\textbf{c，}模型与患者的平均传播秩次。两类模型事件的平均秩次（实线）分别与患者 TA、TB 模板（虚线）总体一致。
\textbf{d，}两类事件的二维传播过程。MTA 和 MTB 分别从 Core A、Core B 附近起始并向相反方向扩展，各列显示距事件起始 0、16、32、64 和 96\,ms 的活动；颜色为按各事件时空峰值归一化的放电率，标尺为 5\,mm。
\textbf{e，}模型传播模式与患者模板的交叉匹配。通过交替触点分折分配事件并评估 Spearman 相关，MTA--TA 和 MTB--TB 呈正相关，交叉组合呈负相关。
\textbf{f，}患者队列中的单、双 core 比较。在 17 位患者的当前搜索结果中，双 core 配置对 TB 的匹配优于偏 TA 的单 core 配置；空心和实心点分别表示单、双 core，细线连接同一患者，横线及竖段表示中位数和四分位范围。
e 采用单侧杆内触点置换检验（1,000 次，未校正）；f 采用配对双侧 Wilcoxon 符号秩检验，并作 Holm 校正。
$^{*}P<0.05$，$^{**}P<0.01$。}
\label{fig5}
\end{figure}
